\documentclass[10pt,a4paper]{amsart}
\usepackage[margin=19mm]{geometry}
\usepackage{amsmath,amssymb,mathtools}
\usepackage[scr=rsfs,cal=euler]{mathalfa}
\usepackage{xfrac}
\usepackage[unicode,hidelinks,bookmarks=false]{hyperref}
\newcommand{\ie}{i.e.\ }
\newcommand{\cf}{cf.\ }
\newcommand{\resp}{respectively\ }
\newcommand{\Subsection}{\subsection}
\providecommand{\nobreakdash}{\nobreak\hbox{-}}
\setlength{\emergencystretch}{2em}
\multlinegap0pt
\usepackage{amssymb}
\usepackage{mathtools}
\newtheorem{lem}{Lemma}
\newcommand{\dif}{\mathrm{d}}
\title{Compactness and Willmore Energy of Helicoidal Minimal Surfaces in the 3-Sphere}
\author{Jianquan Ge, Shilin Li$^{*}$}
\date{Source: arXiv:2607.27965v1 (CC BY 4.0)}
\begin{document}
\maketitle

\noindent\emph{Source and licence.} Compactness and Willmore Energy of Helicoidal Minimal Surfaces in the 3-Sphere. Version arXiv:2607.27965v1 (CC BY 4.0), distributed by arXiv under the Creative Commons Attribution 4.0 International licence, http://creativecommons.org/licenses/by/4.0/, as recorded in the arXiv licence metadata for this exact version and shown on the abstract page https://arxiv.org/abs/2607.27965v1. Full licence notice: \texttt{licenses/arxiv-2607.27965-CC-BY-4.0.txt}.\par\smallskip

\noindent\textit{Editorial scope.} Counted unit: the article's lem lem:q\_range (source0.tex lines 1908--1932). The article's complete original proof of it is delivered with it (source0.tex lines 1934--2110, one \texttt{proof} environment of 177 lines). No mathematics is written, repaired or re-derived: the statement and the proof are the article's own lines, in the article's own notation, and no retained line is substituted or re-flowed.  Retained uncounted as the source-local context the proof needs: lines 989--993; lines 1899--1906.  Nothing was omitted from any retained span, so the package carries no omission record.  No retained line contains a citation, so the wrapper carries no bibliography and no bibliography entry.  No retained line contains a figure environment, an image-inclusion command or drawing code, so no image is delivered and the package declares no asset dependency.  Editorial formatting only, in addition: an AMS \texttt{amsart} preamble that restates the article's own declarations and macros used by the retained lines, namely the environments \texttt{lem} on separate counters, together with the standard packages those lines need.  The retained environments print in this document as Lemma 1 in order of appearance, whereas the article prints the same items with the numbering of its own full document.  The complete native source of the article as distributed in the arXiv e-print for this exact version is bundled unchanged as \texttt{sources/206378/source0.tex}.
\begin{equation}\label{eq:q_definition}
q(h,c)
:=
\frac{\lambda_c^h(s+L)-\lambda_c^h(s)}{2\pi}.
\end{equation}

\begin{equation}\label{eq:q_explicit_integral}
q(h_0,c)
=
\frac{c}{\pi}
\int_{\sqrt{(1-\sqrt{1-4c^2})/2}}^{\sqrt{(1+\sqrt{1-4c^2})/2}}
\frac{\sqrt{h_0^2+(1-h_0^2)z^2}}
{z(1-z^2)\sqrt{z^2-z^4-c^2}}\,\dif z.
\end{equation}

\begin{lem}\label{lem:q_range}
Let $q(h,c)$ be the normalized longitude increment defined by
\eqref{eq:q_definition}. For every fixed $h_0\geq0$, the function
$c\mapsto q(h_0,c)$ is continuous on $(0,1/2)$. Moreover,
\begin{equation}\label{eq:q_endpoint_limits}
\lim_{c\to0^+}q(h_0,c)=\frac{1+h_0}{2},
\qquad
\lim_{c\to1/2^-}q(h_0,c)=\frac{\sqrt{2(1+h_0^2)}}{2}.
\end{equation}
If $h_0\neq1$, then the range of $q(h_0,\cdot)$ on $(0,1/2)$ is the open
interval
\begin{equation}\label{eq:q_range_interval}
\bigl\{q(h_0,c):0<c<1/2\bigr\}
=
\left(
\frac{1+h_0}{2},
\frac{\sqrt{2(1+h_0^2)}}{2}
\right).
\end{equation}
If $h_0=1$, then
\begin{equation}\label{eq:q_pitch_one}
q(1,c)\equiv1
\end{equation}
on $(0,1/2)$, and the range is the singleton $\{1\}$.
\end{lem}

\begin{proof}
We first rewrite $q(h_0,c)$ as an integral over a fixed interval. This
form gives the two endpoint limits by dominated convergence and also yields
strict pointwise bounds between the limiting values.

Set
\[
\delta=\sqrt{1-4c^2},\qquad
A_0=\frac{1+h_0^2}{2},\qquad
B_0=\frac{1-h_0^2}{2}.
\]
Put $t=z^2$, and then from $z^2-z^4-c^2\geq0$ in \eqref{eq:q_explicit_integral} we change the variables by
\[
t=\frac{1-\delta\cos\theta}{2},\qquad 0\le\theta\le\pi.
\]
The elementary identities needed for this change of variables are
\[
    t(1-t)=\frac{1-\delta^2\cos^2\theta}{4},
    \qquad
    t-t^2-c^2=\frac{\delta^2\sin^2\theta}{4},
    \qquad
    \dif t=\frac{\delta}{2}\sin\theta\,\dif\theta.
\]
Moreover
\[
    h_0^2+(1-h_0^2)t=A_0-B_0\delta\cos\theta.
\]
Since
\[
    \frac{\dif z}{z(1-z^2)\sqrt{z^2-z^4-c^2}}
    =
    \frac{\dif t}{2t(1-t)\sqrt{t-t^2-c^2}},
\]
the integral \eqref{eq:q_explicit_integral}  becomes
\[
q(h_0,c)=
\frac{2c}{\pi}
\int_0^\pi
\frac{\sqrt{A_0-B_0\delta\cos\theta}}
{1-\delta^2\cos^2\theta}
\,\dif\theta.
\]
Using the symmetry $\cos(\pi-\theta)=-\cos\theta$, this is equivalently
\begin{equation}\label{eq:q_std}
q(h_0,c)
=
\frac{2c}{\pi}
\int_0^{\pi/2}
\frac{
\sqrt{A_0-B_0\delta\cos\theta}
+
\sqrt{A_0+B_0\delta\cos\theta}
}
{1-\delta^2\cos^2\theta}
\,\dif\theta.
\end{equation}
On every compact subinterval of $(0,1/2)$, the denominator in
\eqref{eq:q_std} is bounded away from zero and the integrand depends
continuously on $c$. Hence $q(h_0,c)$ is continuous on $(0,1/2)$.

Direct calculations show that
\begin{equation}\label{eq:int_aux}
\int_0^{\pi/2}
\frac{\dif\theta}{1-\delta^2\cos^2\theta}
=
\frac{\pi}{4c}.
\end{equation}


For the limit $c\to0^+$, write $u=\tan\theta$ in
\eqref{eq:q_std} and then set $u=2cv$. This gives
\[
q(h_0,c)
=
\frac1\pi
\int_0^\infty
\frac{
\Sigma(\delta,\arctan(2cv))
}
{1+v^2}\,\dif v,
\]
where
\[
\Sigma(\delta,\theta)
=
\sqrt{A_0-B_0\delta\cos\theta}
+
\sqrt{A_0+B_0\delta\cos\theta}.
\]
As $c\to0^+$, $\delta\to1$ and
\[
\Sigma(\delta,\arctan(2cv))\to \Sigma(1,0)=1+h_0.
\]
Since $\Sigma$ is uniformly bounded, dominated convergence gives
\[
\lim_{c\to0^+}q(h_0,c)
=
\frac1\pi
\int_0^\infty
\frac{1+h_0}{1+v^2}\,\dif v
=
\frac{1+h_0}{2}.
\]

For the limit $c\to1/2^-$, we have $\delta\to0$ and $2c\to1$.
For $\delta\leq1/2$, the denominator in \eqref{eq:q_std} is bounded
below by $3/4$, while the numerator is uniformly bounded. Hence dominated
convergence gives
\[
\lim_{c\to1/2^-}q(h_0,c)
=
\frac1\pi
\int_0^{\pi/2}2\sqrt{A_0}\,\dif\theta
=
\sqrt{A_0}
=
\frac{\sqrt{2(1+h_0^2)}}{2}.
\]

If $h_0=1$, then $A_0=1$ and $B_0=0$, so $\Sigma(\delta,\theta)=2$. Using
\eqref{eq:q_std} and \eqref{eq:int_aux}, we get
\[
q(1,c)=1
\]
for every $0<c<1/2$.

Now assume $h_0\neq1$. Then $B_0\neq0$. We verify the strict pointwise
inequality separately because the sign of $B_0$ changes at $h_0=1$.  For
$0\le y\le |B_0|$, set
\[
\mathcal R(y)=\sqrt{A_0-y}+\sqrt{A_0+y}.
\]
Then
\[
\mathcal R'(y)=\frac{1}{2\sqrt{A_0+y}}-\frac{1}{2\sqrt{A_0-y}}<0
\qquad(0<y<|B_0|),
\]
and
\[
\mathcal R(|B_0|)=1+h_0,
\qquad
\mathcal R(0)=2\sqrt{A_0}.
\]
Thus, whenever $0<|B_0|x<|B_0|$,
\[
1+h_0<\mathcal R(|B_0|x)<2\sqrt{A_0}.
\]
Since
\[
\sqrt{A_0-B_0x}+\sqrt{A_0+B_0x}=\mathcal R(|B_0|x),
\]
this applies to $x=\delta\cos\theta$.  For $0<c<1/2$, we have $0<\delta<1$, and thus the strict
inequalities above hold for all $\theta\in[0,\pi/2]$ except for
$\theta=\pi/2$, which is irrelevant for the integral. Combining this with
\eqref{eq:q_std} and
\eqref{eq:int_aux}, we obtain the strict bounds
\[
\frac{1+h_0}{2}
<
q(h_0,c)
<
\sqrt{A_0}
=
\frac{\sqrt{2(1+h_0^2)}}{2},\quad \text{for}\quad 0<c<1/2.
\]
Since $q(h_0,\cdot)$ is continuous, its image is an interval.  The endpoint limits in \eqref{eq:q_endpoint_limits} give its infimum and
supremum, while the strict inequalities exclude the endpoints. Therefore, for $h_0\neq1$,
\[
q(h_0,(0,1/2))
=
\left(
\frac{1+h_0}{2},
\frac{\sqrt{2(1+h_0^2)}}{2}
\right).
\]
This proves the lemma.
\end{proof}

\end{document}
